\documentclass[../Main.tex]{subfiles}
\begin{document}

\section{Conclusion}
\label{section:6.1}

This graduation thesis has designed, implemented, and evaluated a \textbf{Smart Medical QA and Consultation System} combining Graph Retrieval-Augmented Generation (GraphRAG) and Agentic AI. The system operates entirely offline on consumer hardware, ensuring full patient data confidentiality.

The key achievements of this thesis are summarized as follows:

\begin{itemize}
    \item \textbf{Purified Clinical Knowledge Graph}: A custom medical KG was constructed from 9,989 clinical documents. After systematic entity pruning to remove conversational noise, the graph retains 25,319 high-quality clinical entities and 193,042 semantic relations, indexed in Neo4j for direct Cypher-based retrieval.
    \item \textbf{Graph-First Hybrid Retrieval}: A dual-layer retrieval engine was implemented, routing queries first to the Custom KG for precise local entity matching, and falling back to Microsoft GraphRAG community summaries for global synthesis when needed.
    \item \textbf{Resilient ReAct Agent}: A robust agentic reasoning loop was built on top of a local Ollama LLM, reinforced with three defensive mechanisms — parse-retry prompting, iteration-1 recovery-forced GraphRAG, and loop-guard signature tracking — achieving a 100\% meaningful-response rate on adversarial test queries.
    \item \textbf{Multi-Format Clinical Analyzer}: An automatic document processing pipeline supports PDF and Excel medical records, performs on-form reference range benchmarking without hardcoded tables, and generates graph-grounded clinical consultations enriched with pill imagery.
    \item \textbf{Triangulated Reranking and Semantic Graph Pruning}: Fused dual-channel retrieval results via Reciprocal Rank Fusion (RRF) and integrated a neural Cross-Encoder model locally to boost Precision@5 to 0.245, combined with a graph soft-pruning algorithm to systematically sweep out isolated/noisy Neo4j nodes.
    \item \textbf{CDSS Workflows and Conversational Memory}: Closed the E2E clinical workflow by implementing a local conversational chat memory with regex-based HTML sanitization and a high-fidelity vector clinical PDF report exporter, transforming a standard chat application into a fully-featured Clinical Decision Support System (CDSS).
    \item \textbf{Production-Ready Deployment}: The entire system is containerized using Docker Compose, covering the API server, Neo4j graph database, and Ollama LLM runner, fully reproducible from a single command.
\end{itemize}

Compared to baseline approaches, the proposed system demonstrates clear advantages over standard Vector RAG (no structural cross-referencing) and commercial cloud-based LLMs (privacy violation, API cost). Its Graph-First architecture provides explicit, traceable evidence citations for every answer, a property absent from dense retrieval systems.

During the development process, several important engineering lessons were learned: the importance of aggressive entity pruning in clinical KG construction; the necessity of defensive parsing in small LLM deployments; the value of separating pure-Python extraction logic from LLM synthesis to maximize reliability; and the importance of offline-first design when handling sensitive health data.

\section{Future Work}
\label{section:6.2}

While the current system achieves its primary objectives, several directions exist for further improvement:

\textbf{Short-term improvements:}
\begin{itemize}
    \item \textbf{Semantic Caching}: Implement a local vector-similarity cache layer (e.g., lightweight local cosine embeddings) before the ReAct loop to detect semantically identical medical queries (similarity > 0.95), serving historical consultations instantly (under 5ms) and further decreasing token overhead.
    \item \textbf{Expanded evaluation dataset}: The current evaluation uses only 10 clinical queries. Expanding the dataset to 100+ queries covering a broader range of disease categories, drug interactions, and lab indicators would provide statistically significant benchmark results.
\end{itemize}

\textbf{Medium-term improvements:}
\begin{itemize}
    \item \textbf{Asynchronous document processing queue}: The medical record analysis pipeline currently runs synchronously in the FastAPI request thread. For large Excel files or multi-page PDFs, this can cause request timeouts. Migrating to an asynchronous task queue (e.g., Celery with Redis) would allow background processing with webhook callbacks or WebSocket progress updates.
    \item \textbf{LLM model upgrade}: As larger, more instruction-following models (e.g., Llama-3.3-70B via quantization, or Qwen2.5-14B) become feasible on consumer hardware, upgrading the local LLM backend is expected to substantially reduce ReAct format errors and improve answer quality.
\end{itemize}

\textbf{Long-term improvements:}
\begin{itemize}
    \item \textbf{Full observability stack}: Integrating distributed tracing (OpenTelemetry), Prometheus metrics for per-endpoint latency monitoring, and structured logging for agent reasoning traces would enable production-grade monitoring and anomaly detection.
    \item \textbf{Federated KG with hospital systems}: The current KG is built from a single public conversational medical dataset. Federated integration with hospital FHIR APIs or clinical trial databases would expand knowledge coverage without centralizing sensitive patient data.
    \item \textbf{Multimodal clinical inputs}: Extending the document analyzer to support radiology imaging reports (DICOM), ECG trace analysis, and handwritten prescription OCR would broaden the system's applicability across clinical specialties.
\end{itemize}

\end{document}